% GENERATED FILE. Edit canonical lessons, then run npm run generate:books.
% canonical-source-hash: fnv1a64:1d7274cfb3962d5d
% canonical-lessons: ML-C425-pratisedham, ML-C425-pratisedhikkuka, ML-C425-samaram, ML-C425-harttal, ML-C425-jatha

\chapter{Protest}
\label{ch:ml-protest}
\hlchaptermodality{425}

\begin{chapteropening}
\textbf{By the end of this chapter:} \emph{I can say a protest, to protest, a strike, a general strike and a march.}

Complete the last lesson of chapter 425: I can say a protest, to protest, a strike, a general strike and a march.
\end{chapteropening}

\section[\texorpdfstring{\ml{പ്രതിഷേധം} (pratiṣēdhaṁ)}{പ്രതിഷേധം (pratiṣēdhaṁ)}]{\ml{പ്രതിഷേധം} (pratiṣēdhaṁ) — a protest}

\label{lesson:ML-C425-pratisedham}



\begin{quote}
Before the new one: say the Malayalam for authority, then the Malayalam for a leader.
\end{quote}

\subsection*{You'll want to know: \ml{പ്രതിഷേധം}}
\textbf{\ml{പ്രതിഷേധം}} — \emph{pratiṣēdhaṁ} — \textquotedblleft{}a protest\textquotedblright{}.

\begin{quote}
\ml{വില} \ml{കൂടിയതിനെതിരെ} \ml{നഗരത്തിൽ} \ml{പ്രതിഷേധം} \ml{നടന്നു}.
\end{quote}

\begin{quote}
\emph{vila kūṭiyatinetire nagarattil pratiṣēdhaṁ naṭannu.}
\end{quote}

\begin{quote}
There was a protest in the city against the price rise.
\end{quote}

Against something: -\ml{ന്} \ml{എതിരെ} \ml{പ്രതിഷേധം}.

\begin{cousinweb}
\textbf{\ml{പ്രതിഷേധം}} (\emph{pratiṣēdhaṁ}) is borrowed from Sanskrit \emph{pratiṣedha}, \textquotedblleft{}keeping back, warding off, prohibition\textquotedblright{} (Monier-Williams). A protest pushes back against a decision.
\end{cousinweb}

\begin{grammarlens}[title={What you've built}]
One more word for giving reasons and opinions.
\end{grammarlens}

\subsection*{Your turn}
\begin{itemize}
\raggedright
  \item \emph{Say it:} \emph{pratiṣēdhaṁ}
  \item \emph{Say it:} \emph{pratiṣēdhaṁ}, once more
  \item \emph{Say it:} \emph{vila kūṭiyatinetire nagarattil pratiṣēdhaṁ naṭannu.} — There was a protest in the city against the price rise.
  \item \emph{Recall:} say the Malayalam for authority, then the Malayalam for a leader, then say \emph{pratiṣēdhaṁ} again
\end{itemize}

\begin{tcolorbox}[breakable,colback=teal!4,colframe=teal!35!black,title={Before you move on}]
What does \ml{പ്രതിഷേധം} mean? (\textquotedblleft{}A protest\textquotedblright{}.) Say it once more.
\end{tcolorbox}
\section[\texorpdfstring{\ml{പ്രതിഷേധിക്കുക} (pratiṣēdhikkuka)}{പ്രതിഷേധിക്കുക (pratiṣēdhikkuka)}]{\ml{പ്രതിഷേധിക്കുക} (pratiṣēdhikkuka) — to protest}

\label{lesson:ML-C425-pratisedhikkuka}



\begin{quote}
Before the new one: say the Malayalam for a leader, then the Malayalam for a protest.
\end{quote}

\subsection*{You'll want to know: \ml{പ്രതിഷേധിക്കുക}}
\textbf{\ml{പ്രതിഷേധിക്കുക}} — \emph{pratiṣēdhikkuka} — \textquotedblleft{}to protest\textquotedblright{}.

\begin{quote}
\ml{ഫീസ്} \ml{കൂട്ടിയതിൽ} \ml{വിദ്യാർത്ഥികൾ} \ml{പ്രതിഷേധിച്ചു}.
\end{quote}

\begin{quote}
\emph{phīsŭ kūṭṭiyatil vidyārtthikaḷ pratiṣēdhiccu.}
\end{quote}

\begin{quote}
The students protested against the fee increase.
\end{quote}

Past: \ml{പ്രതിഷേധിച്ചു}, \emph{pratiṣēdhiccu}.

The cause takes -\ml{ഇൽ}: \ml{ഫീസ്} \ml{കൂട്ടിയതിൽ}, over raising the fee.

\begin{cousinweb}
\textbf{\ml{പ്രതിഷേധിക്കുക}} (\emph{pratiṣēdhikkuka}) is built from Sanskrit \emph{pratiṣedha} (Monier-Williams), with the Malayalam verb ending -\ml{ിക്കുക}. It is the verb of \ml{പ്രതിഷേധം}.
\end{cousinweb}

\begin{grammarlens}[title={What you've built}]
One more word for giving reasons and opinions.
\end{grammarlens}

\subsection*{Your turn}
\begin{itemize}
\raggedright
  \item \emph{Say it:} \emph{pratiṣēdhikkuka}
  \item \emph{Say it:} \emph{pratiṣēdhikkuka}, once more
  \item \emph{Say it:} \emph{phīsŭ kūṭṭiyatil vidyārtthikaḷ pratiṣēdhiccu.} — The students protested against the fee increase.
  \item \emph{Recall:} say the Malayalam for a leader, then the Malayalam for a protest, then say \emph{pratiṣēdhikkuka} again
\end{itemize}

\begin{tcolorbox}[breakable,colback=teal!4,colframe=teal!35!black,title={Before you move on}]
What does \ml{പ്രതിഷേധിക്കുക} mean? (\textquotedblleft{}To protest\textquotedblright{}.) Say it once more.
\end{tcolorbox}
\section[\texorpdfstring{\ml{സമരം} (samaraṁ)}{സമരം (samaraṁ)}]{\ml{സമരം} (samaraṁ) — a strike; a struggle}

\label{lesson:ML-C425-samaram}



\begin{quote}
Before the new one: say the Malayalam for a protest, then the Malayalam for to protest.
\end{quote}

\subsection*{You'll want to know: \ml{സമരം}}
\textbf{\ml{സമരം}} — \emph{samaraṁ} — \textquotedblleft{}a strike; a struggle\textquotedblright{}.

\begin{quote}
\ml{തൊഴിലാളികൾ} \ml{ഇന്ന്} \ml{മുതൽ} \ml{സമരം} \ml{തുടങ്ങി}.
\end{quote}

\begin{quote}
\emph{toḻilāḷikaḷ innŭ mutal samaraṁ tuṭaṅṅi.}
\end{quote}

\begin{quote}
The workers began a strike from today.
\end{quote}

To strike is \ml{സമരം} \ml{ചെയ്യുക}.

\begin{cousinweb}
\textbf{\ml{സമരം}} (\emph{samaraṁ}) is borrowed from Sanskrit \emph{samara}, \textquotedblleft{}a coming together; a hostile encounter, a struggle\textquotedblright{} (Monier-Williams).
\end{cousinweb}

\begin{grammarlens}[title={What you've built}]
One more word for giving reasons and opinions.
\end{grammarlens}

\subsection*{Your turn}
\begin{itemize}
\raggedright
  \item \emph{Say it:} \emph{samaraṁ}
  \item \emph{Say it:} \emph{samaraṁ}, once more
  \item \emph{Say it:} \emph{toḻilāḷikaḷ innŭ mutal samaraṁ tuṭaṅṅi.} — The workers began a strike from today.
  \item \emph{Recall:} say the Malayalam for a protest, then the Malayalam for to protest, then say \emph{samaraṁ} again
\end{itemize}

\begin{tcolorbox}[breakable,colback=teal!4,colframe=teal!35!black,title={Before you move on}]
What does \ml{സമരം} mean? (\textquotedblleft{}A strike; a struggle\textquotedblright{}.) Say it once more.
\end{tcolorbox}
\section[\texorpdfstring{\ml{ഹർത്താൽ} (harttāl)}{ഹർത്താൽ (harttāl)}]{\ml{ഹർത്താൽ} (harttāl) — a general strike}

\label{lesson:ML-C425-harttal}



\begin{quote}
Before the new one: say the Malayalam for to protest, then the Malayalam for a strike.
\end{quote}

\subsection*{You'll want to know: \ml{ഹർത്താൽ}}
\textbf{\ml{ഹർത്താൽ}} — \emph{harttāl} — \textquotedblleft{}a general strike\textquotedblright{}.

\begin{quote}
\ml{ഹർത്താൽ} \ml{കാരണം} \ml{കടകൾ} \ml{എല്ലാം} \ml{അടഞ്ഞുകിടന്നു}.
\end{quote}

\begin{quote}
\emph{harttāl kāraṇaṁ kaṭakaḷ ellāṁ aṭaññukiṭannu.}
\end{quote}

\begin{quote}
Because of the general strike, all the shops stayed shut.
\end{quote}

On a \ml{ഹർത്താൽ} day, shops close and buses stop.

\begin{cousinweb}
\textbf{\ml{ഹർത്താൽ}} (\emph{harttāl}) was borrowed from Hindustani \emph{haṛtāl}, \textquotedblleft{}closing shops in mourning or protest\textquotedblright{}, which CDIAL lists under Sanskrit \emph{haṭṭa}, \textquotedblleft{}a market\textquotedblright{} (CDIAL 13944).
\end{cousinweb}

\begin{grammarlens}[title={What you've built}]
One more word for giving reasons and opinions.
\end{grammarlens}

\subsection*{Your turn}
\begin{itemize}
\raggedright
  \item \emph{Say it:} \emph{harttāl}
  \item \emph{Say it:} \emph{harttāl}, once more
  \item \emph{Say it:} \emph{harttāl kāraṇaṁ kaṭakaḷ ellāṁ aṭaññukiṭannu.} — Because of the general strike, all the shops stayed shut.
  \item \emph{Recall:} say the Malayalam for to protest, then the Malayalam for a strike, then say \emph{harttāl} again
\end{itemize}

\begin{tcolorbox}[breakable,colback=teal!4,colframe=teal!35!black,title={Before you move on}]
What does \ml{ഹർത്താൽ} mean? (\textquotedblleft{}A general strike\textquotedblright{}.) Say it once more.
\end{tcolorbox}
\section[\texorpdfstring{\ml{ജാഥ} (jātha)}{ജാഥ (jātha)}]{\ml{ജാഥ} (jātha) — a march; a procession}

\label{lesson:ML-C425-jatha}



\begin{quote}
Before the new one: say the Malayalam for a strike, then the Malayalam for a general strike.
\end{quote}

\subsection*{You'll want to know: \ml{ജാഥ}}
\textbf{\ml{ജാഥ}} — \emph{jātha} — \textquotedblleft{}a march; a procession\textquotedblright{}.

\begin{quote}
\ml{കുട്ടികൾ} \ml{ഗ്രാമത്തിലൂടെ} \ml{ഒരു} \ml{ജാഥ} \ml{നടത്തി}.
\end{quote}

\begin{quote}
\emph{kuṭṭikaḷ grāmattilūṭe oru jātha naṭatti.}
\end{quote}

\begin{quote}
The children held a march through the village.
\end{quote}

A march is \ml{നടത്തുക}, organised.

\begin{cousinweb}
\textbf{\ml{ജാഥ}} (\emph{jātha}) is usually traced to Hindustani \emph{jathā}, \textquotedblleft{}a band, a group of people\textquotedblright{}; its history before that is not settled.
\end{cousinweb}

\begin{grammarlens}[title={What you've built}]
One more word for giving reasons and opinions.
\end{grammarlens}

\subsection*{Your turn}
\begin{itemize}
\raggedright
  \item \emph{Say it:} \emph{jātha}
  \item \emph{Say it:} \emph{jātha}, once more
  \item \emph{Say it:} \emph{kuṭṭikaḷ grāmattilūṭe oru jātha naṭatti.} — The children held a march through the village.
  \item \emph{Recall:} say the Malayalam for a strike, then the Malayalam for a general strike, then say \emph{jātha} again
\end{itemize}

\begin{tcolorbox}[breakable,colback=teal!4,colframe=teal!35!black,title={Before you move on}]
What does \ml{ജാഥ} mean? (\textquotedblleft{}A march; a procession\textquotedblright{}.) Say it once more.
\end{tcolorbox}
